\documentclass[11pt,a4paper]{article}
\usepackage[T1]{fontenc}
\usepackage[utf8]{inputenc}
\usepackage[czech]{babel}
\usepackage{lmodern,microtype}
\usepackage{amsmath,amssymb,mathtools}
\usepackage{geometry}
\geometry{left=19mm,right=19mm,top=21mm,bottom=21mm}
\usepackage[dvipsnames]{xcolor}
\usepackage{enumitem,tabularx,booktabs}
\usepackage[most]{tcolorbox}
\usepackage{fancyhdr,lastpage}
\usepackage[hidelinks]{hyperref}
\definecolor{deep}{HTML}{203955}
\definecolor{teal}{HTML}{176E76}
\definecolor{slate}{HTML}{506173}
\definecolor{pale}{HTML}{F3F7F9}
\setlength{\parindent}{0pt}
\setlength{\parskip}{5pt}
\setlist[enumerate]{leftmargin=6mm,itemsep=5pt,topsep=3pt}
\pagestyle{fancy}\fancyhf{}
\fancyhead[L]{\sffamily\footnotesize\color{slate} MAT · 4. ročník · kombinatorika}
\fancyhead[R]{\sffamily\footnotesize\color{slate} Hodina \LESSONNUM\, / 112}
\fancyfoot[L]{\sffamily\footnotesize\color{slate} Výukový podklad · samostatný dokument}
\fancyfoot[R]{\sffamily\footnotesize\color{slate} \thepage/\pageref{LastPage}}
\renewcommand{\headrulewidth}{0.35pt}
\renewcommand{\footrulewidth}{0pt}
\newcommand{\sect}[1]{\par\vspace{2mm}{\Large\sffamily\bfseries\color{deep} #1}\par\vspace{-0.5mm}\noindent\color{teal}\rule{\linewidth}{0.7pt}\color{black}\vspace{1mm}}
\tcbset{theobox/.style={enhanced,breakable,colback=pale,colframe=teal!45!white,boxrule=.6pt,arc=1mm,left=3mm,right=3mm,top=2mm,bottom=2mm}}
\newcommand{\LESSONNUM}{34}
\begin{document}

{\Huge\sffamily\bfseries\color{deep} Variace s opakováním a kódy}\par
{\sffamily\large\color{teal} Výukový podklad · 34. vyučovací hodina}\par
\textbf{Hodinová dotace:} 45 minut\quad \textbf{Ročník:} 4. SPŠ\quad \textbf{Rozpis:} TP 2.1\par
\begin{tcolorbox}[theobox]\textbf{Výukový cíl.} Odvodit \(n^k\), poznat rozdíl mezi kódem a číslem, řešit úlohy s povoleným nebo zakázaným opakováním.\end{tcolorbox}
\sect{Teorie a vysvětlení}
\textbf{Variace s opakováním} délky \(k\) z \(n\) různých symbolů: na každém místě lze vybrat kterýkoli symbol, pořadí rozhoduje. Počet \(\overline V(n,k)=n^k\).\par
\textbf{Proti tomu variace bez opakování}: \(V(n,k)=n(n-1)\cdots(n-k+1)\), neboť po každé volbě ubude jeden dovolený symbol.\par
\textbf{Doplňková metoda.} Počet výsledků obsahujících aspoň jeden zakázaný/daný znak lze často spočítat jako \(\text{všechny}-\text{žádný}\). Je-li zakázán první znak \(0\), musíme samostatně vyřešit první pozici.
\sect{Řešené příklady}
\textbf{Příklad 1.} Kolik čtyřmístných PINů používá číslice \(0\) až \(9\)?\par
\textit{Řešení.} Kód může začít nulou a číslice se mohou opakovat. \(10^4=10\,000\).\par\smallskip
\textbf{Příklad 2.} Kolik je čtyřciferných přirozených čísel při dovoleném opakování?\par
\textit{Řešení.} První pozice \(9\) možností (\(1\)--\(9\)), zbývající tři \(10\). \(9\cdot10^3=9000\), nikoli \(10^4\).\par\smallskip
\textbf{Příklad 3.} Kolik trojznakových kódů sestavíme z šesti písmen, bez opakování?\par
\textit{Řešení.} Máme \(6\) možností na první pozici, \(5\) na druhou a \(4\) na třetí. \(6\cdot5\cdot4=120\). Pro porovnání s opakováním by to bylo \(6^3=216\).\par\smallskip
\textbf{Příklad 4.} Trojznakové heslo z \(A,B,C,D\). Kolik hesel obsahuje právě jedno \(A\)?\par
\textit{Řešení.} Vybereme pozici znaku \(A\) (\(3\) možnosti); zbývající dvě pozice obsadíme nezávisle jedním z \(B,C,D\) (\(3^2\) možností). \(3\cdot3^2=27\).\par\smallskip
\textbf{Příklad 5.} Kolik trojciferných kódů z \(0\) až \(9\) obsahuje alespoň jednu sedmičku?\par
\textit{Řešení.} Ve všech kódech je \(10^3\) možností. Žádná sedmička: \(9^3\). Alespoň jedna sedmička: \(1000-729=271\).\par\smallskip
\clearpage\sect{Procvičování pro studenty}
\begin{enumerate}
\item Kolik pětic symbolů vznikne ze tří znaků, s opakováním?
\item Kolik dvoupísmenných hesel vznikne z pěti různých písmen bez opakování?
\item Kolik čtyřznakových kódů z pěti symbolů neobsahuje symbol X?
\item Kolik čtyřznakových kódů ze tří znaků obsahuje alespoň jedno \(A\)?
\item Kolik čtyřmístných čísel lze vytvořit z cifer \(0\) až \(5\), s opakováním?
\end{enumerate}
\sect{Řešení a zdůvodnění k procvičování}
\begin{enumerate}
\item \(3^5=243\).
\item \(5\cdot4=20\).
\item Na každé pozici zbývají \(4\) možnosti: \(4^4=256\).
\item \(3^4-2^4=81-16=65\).
\item První cifra \(5\) možností, ostatní \(6\): \(5\cdot6^3=1080\).
\end{enumerate}
\sect{Metodické poznámky}
\textbf{Typické chyby:}
\begin{itemize}[leftmargin=5mm,itemsep=1pt]
\item \(n^k\) neplatí pro výběr bez opakování.
\item Heslo/pin může začínat nulou, číslo ne.
\item „Alespoň jedno“ zpravidla počítat doplňkem, ne výčtem mnoha překrývajících se případů.
\end{itemize}
\textbf{Doporučený průběh 45 minut:} 0–8 porovnání PINu a čísla; 8–18 \(n^k\) a \(V(n,k)\); 18–30 omezení; 30–42 úlohy; 42–45 kontrola.\par
\textbf{Návaznost:} Příště na pořadí přestane záležet a budeme dělit počet pořadí téhož výběru.\par
\vspace{2mm}\noindent{\footnotesize\color{slate} \textbf{Východiska:} rozpis výuky MAT, 4. ročník, hodiny 29--41; tematický plán SPŠ; srovnání tematického rozsahu a náročnosti s realisticky.cz, oddíl Kombinatorika (lekce 090101--090119). Výklad, příklady a řešení jsou nově autorsky sestavené, nikoli převzatá sbírka.}\par
\end{document}